\documentclass[11pt]{article}

\usepackage[margin=1in]{geometry}
\usepackage{amsmath,amssymb}
\usepackage{booktabs}
\usepackage{microtype}
\usepackage{natbib}
\usepackage{setspace}
\usepackage{threeparttable}
\usepackage{xcolor}
\usepackage[colorlinks=true,allcolors=blue]{hyperref}

\setlength{\parindent}{1.5em}
\setlength{\parskip}{0pt}
\onehalfspacing

\title{Resolution Risk in Prediction Markets}
\author{Terry Cheng}
\date{June 2026}

\begin{document}
\maketitle

\begin{abstract}
Prediction-market prices are commonly interpreted as probabilities of underlying events.
That interpretation assumes that the exchange will map the realized event into the intended
contract payoff. This project studies the possibility that prices also reflect
\emph{resolution risk}: uncertainty that a contract will be resolved in a manner that differs
from traders' economic interpretation of its terms. We construct high-precision matches
between Polymarket and Kalshi contracts and study two 2025 controversies involving
Polymarket's UMA-based resolution process. The primary outcomes measure each price's
absolute distance from 0.5 and Polymarket's distance relative to Kalshi's. Preliminary
estimates do not support the simple hypothesis that a public resolution controversy causes
prices broadly to attenuate toward 0.5. The exercise is informative, but it is not yet a
causal estimate: the available data, matching procedure, event windows, weights, and standard
errors all require further work.
\end{abstract}

\section{Introduction}

A binary prediction-market price is often read as the market's probability that an event
will occur. Even in an ideal contract this interpretation requires assumptions about beliefs,
risk preferences, wealth, and trading frictions \citep{manski2006,wolfers2006}. In practice,
it also requires confidence in the institution that converts the realized event into a
payoff. A trader may believe that an event occurred while assigning positive probability to
the exchange resolving the contract the other way. The contract price then reflects both
beliefs about the event and beliefs about resolution.

This project studies that second component, which we call \emph{resolution risk}. The
institutional comparison is between Polymarket and Kalshi. During the period studied here,
the matched Polymarket contracts were ultimately backed by UMA's optimistic-oracle and
dispute process, either directly or through Polymarket's negative-risk adapter. Kalshi
resolves contracts internally: the exchange applies its own contract rules and settlement
procedures rather than relying on UMA or another tokenholder-based dispute mechanism.
Equivalent contracts on the two exchanges are therefore exposed to similar event news but
different resolution institutions.

We focus on two public controversies: the March 24, 2025 Ukraine mineral-rights resolution
and the June 30--July 8, 2025 Zelensky-suit episode. The research question is whether these
episodes changed the informativeness of Polymarket prices relative to matched Kalshi prices.
Our main hypothesis is attenuation. If a controversy adds uncertainty about which side of a
contract will be paid, then a Polymarket price should move toward 0.5 relative to an otherwise
equivalent Kalshi price. This prediction is deliberately symmetric: the relevant object is
not whether the average Yes price rises or falls, but whether Polymarket prices become less
extreme.

\section{Related Literature and Contribution}

The project is closest to three literatures. First, prediction-market prices need not equal
literal probabilities, even when markets aggregate dispersed information
\citep{wolfers2004,manski2006,wolfers2006}. Resolution risk supplies a distinct institutional
wedge: a claim can be priced differently even when traders agree about the underlying event
because they disagree about how the contract will be interpreted or enforced.

Second, work on manipulation in prediction markets asks whether strategic trading can
distort prices or whether traders can change the event being forecast
\citep{hanson2006,ottaviani2007,rhode2008}. Our mechanism is different. The object at risk is
the mapping from the event to the payoff, not necessarily the event itself or the trading
price before resolution. This connects the paper to recent work on semantic ambiguity in
prediction-market contracts \citep{fisman2026} and to the computer-science literature on
blockchain oracles, governance, and manipulation-resistant data feeds
\citep{eskandari2021}.

Third, matched contracts across exchanges resemble law-of-one-price comparisons. Limits to
arbitrage, collateral constraints, segmentation, and market-specific liquidity can sustain
price wedges even for closely related claims \citep{shleifer1997,froot1999,garleanu2011}.
These forces are not merely background citations; they are competing explanations that the
empirical design must address. The intended contribution is therefore narrow: to test
whether public shocks to the perceived credibility of a resolution institution alter the
informativeness of otherwise comparable prediction-market prices.

\section{Institutions and Hypothesis}

UMA uses an optimistic-oracle design. An outcome is proposed, a challenge period allows a
dispute, and disputed cases can ultimately reach UMA's decentralized verification mechanism.
This structure is designed to make incorrect resolution costly, but it also makes the final
payoff depend on an institutional process that is distinct from the underlying event.
Polymarket's negative-risk markets add an adapter layer, but the matched sample's ultimate
resolver remains UMA.

Let $p^{P}_{it}$ and $p^{K}_{it}$ denote the Polymarket and Kalshi prices for matched contract
$i$ at time $t$. Define each price's extremity as
\begin{align}
A^{P}_{it} &= \left|p^{P}_{it}-0.5\right|, \\
A^{K}_{it} &= \left|p^{K}_{it}-0.5\right|,
\end{align}
and define relative extremity as
\begin{equation}
R_{it}=A^{P}_{it}-A^{K}_{it}.
\end{equation}
A negative change in $A^{P}_{it}$ indicates that Polymarket moves toward 0.5 after a
controversy. A negative change in $R_{it}$ indicates attenuation on Polymarket relative to
the matched Kalshi contract. The second outcome is the cleaner comparison because it removes
common changes in how close the event itself is to resolution. We do not emphasize a
directional Polymarket--Kalshi price difference because resolution uncertainty can impair
either side of a binary contract; its sign depends on which outcome is vulnerable.

\section{Data and Matching}

\subsection{Source data and price construction}

The analysis uses the J.D. Becker archive of Polymarket and Kalshi market and trade data.
For each event window, we construct observations at 5:00 a.m. and 5:00 p.m. Pacific using
the most recent trade at or before the cutoff. The analysis is therefore a panel of
\emph{as-of price snapshots}, not a synchronized panel of contemporaneous quotes or order
books. A price may be many hours or days old when a market trades infrequently.

This distinction is especially important for nonprice variables. The archive and
pipeline do not provide aligned time series for bid--ask spreads, depth, open interest, or
exchange-level participation. Polymarket volume enters as a market-level snapshot rather
than a contemporaneous, event-time measure, and comparable Kalshi volume histories are not
used. Consequently, the analysis cannot determine whether a price response is
accompanied by a change in liquidity or participation, and its volume weights are ex post.

\subsection{Deterministic matching}

The matching pipeline parses each contract into a structured key containing the market
family and relevant semantic fields, such as predicate, subject, object, threshold, and
deadline. Aliases and economically equivalent deadlines are normalized before exact keys
are joined across exchanges. The parser covers 17 contract families, including
elections, office exits, leader meetings, sports, Federal Reserve decisions, policy actions,
and asset-price thresholds.

The procedure begins with 5,680 Polymarket markets and 8,992 Kalshi markets. It parses 1,346
Polymarket markets and 1,979 Kalshi markets, then produces 150 deterministic matched rows
representing 149 unique keys. Parsing rates are low in counts---23.7 percent for Polymarket
and 22.0 percent for Kalshi---but much higher when weighted by reported volume. On-chain
checks classify 15 matched Polymarket markets as direct UMA markets and 135 as resolving
through the negative-risk adapter to UMA.

A more flexible matcher could improve coverage without sacrificing the accuracy of the
deterministic approach by separating candidate generation from adjudication.
Markets could first be segmented by broad family, date range, and named entities. Within
each segment, embeddings or lexical retrieval could generate a small candidate set. A
language model could then compare the full contract text and resolution rules using a fixed
schema, returning a match decision, the specific rule differences, and a calibrated
confidence score. Human review of accepted matches would preserve a high-precision final
sample while extending the method to contracts not covered by hand-written parsers.

\section{Empirical Design}

For each controversy, the empirical analysis constructs a symmetric seven-day window around the
event period and estimates
\begin{equation}
Y_{it}=\alpha_i+\beta\,\mathrm{Post}_{t}+\varepsilon_{it},
\label{eq:empirical}
\end{equation}
where $Y_{it}$ is either $A^{P}_{it}$ or $R_{it}$ and $\alpha_i$ is a matched-contract fixed
effect. Observations are weighted by the available Polymarket volume measure. The mineral-
rights event contributes 105 matched pairs and 3,105 observations; the Zelensky-suit event
contributes 57 pairs and 2,565 observations. A narrower ``resolution-risk-exposed'' sample
keeps elections, office exits, leader contacts, and policy actions, leaving 12 and 19 pairs
for the two episodes.

Equation \eqref{eq:empirical} is a pre/post fixed-effect comparison, not a modern dynamic
event study or a fully specified difference-in-differences design. Kalshi enters through the
construction of $R_{it}$ rather than through an exchange-level panel with matched-pair-by-
time fixed effects. The present specification also imposes a single discontinuity at the
start of each episode even though the Zelensky controversy spans multiple days.

\section{Preliminary Results}

Table \ref{tab:results} reports the estimates in the conventional economics format:
coefficients followed by standard errors in parentheses. Negative coefficients would support
attenuation toward 0.5. The estimates are generally positive. In the full sample, the
mineral-rights episode produces essentially no relative change, while the Zelensky episode
is associated with a 0.47 percentage point increase in Polymarket's extremity relative to
Kalshi. In the resolution-risk-exposed sample, the corresponding Zelensky estimate is a
1.47 percentage point increase. These patterns reject the simplest version of the
attenuation hypothesis in this specification.

\begin{table}[!htbp]
\centering
\caption{Resolution Controversies and Price Extremity}
\label{tab:results}
\begin{threeparttable}
\small
\begin{tabular*}{\textwidth}{@{\extracolsep{\fill}}lcccc}
\toprule
& \multicolumn{2}{c}{Mineral rights} & \multicolumn{2}{c}{Zelensky suit} \\
\cmidrule(lr){2-3}\cmidrule(lr){4-5}
& $|P-.5|$ & $|P-.5|-|K-.5|$ & $|P-.5|$ & $|P-.5|-|K-.5|$ \\
& (1) & (2) & (3) & (4) \\
\midrule
\multicolumn{5}{l}{\textit{Panel A: All matched contracts}} \\
Post
    & 0.0010 & 0.0003 & 0.0197 & 0.0047 \\
    & (0.0004) & (0.0004) & (0.0014) & (0.0008) \\
Pre-event mean
    & 0.4564 & 0.0071 & 0.3728 & -0.0197 \\
Observations
    & 3,105 & 3,105 & 2,565 & 2,565 \\
Matched contracts
    & 105 & 105 & 57 & 57 \\
\addlinespace
\multicolumn{5}{l}{\textit{Panel B: Resolution-risk-exposed contracts}} \\
Post
    & 0.0117 & 0.0012 & 0.0204 & 0.0147 \\
    & (0.0031) & (0.0027) & (0.0016) & (0.0013) \\
Pre-event mean
    & 0.1969 & 0.0356 & 0.3437 & -0.0596 \\
Observations
    & 355 & 355 & 817 & 817 \\
Matched contracts
    & 12 & 12 & 19 & 19 \\
\midrule
Contract fixed effects & Yes & Yes & Yes & Yes \\
Polymarket-volume weights & Yes & Yes & Yes & Yes \\
\bottomrule
\end{tabular*}
\begin{tablenotes}[flushleft]
\footnotesize
\item \textit{Notes:} Each observation is an as-of price snapshot for a deterministic
Polymarket--Kalshi contract match. The dependent variables are Polymarket's absolute distance
from 0.5 and that distance minus the matched Kalshi contract's absolute distance from 0.5.
The reported coefficient is from a weighted regression of the dependent variable on a post
indicator and matched-contract fixed effects. Parentheses contain the conventional standard
errors used in the empirical analysis. They are not clustered and should not be used for
formal inference. No significance stars are reported. The resolution-risk-exposed sample
contains elections, office exits, leader contacts, and policy actions.
\end{tablenotes}
\end{threeparttable}
\end{table}

The results do not establish that resolution concerns had no effect. The hypothesized
response may be concentrated in the disputed contracts themselves, in contracts near their
terminal dates, or in markets whose rules contain genuine interpretive ambiguity. It may
also appear in spreads, depth, withdrawal of liquidity, or trading activity rather than in
price attenuation. Conversely, the positive coefficients may simply reflect stale prices,
ordinary event news, compositional changes, or the mechanical behavior of absolute-value
outcomes. The present estimates cannot distinguish these explanations.

\section{Limitations and Required Empirical Work}

The analysis has four central limitations.

\paragraph{Data frequency and market quality.}
Last-trade snapshots are not synchronized executable prices. In the constructed panels,
Kalshi observations are substantially staler than Polymarket observations, with long tails
of contracts that have not traded for days. The analysis should impose maximum trade-age
restrictions and, where available, use synchronized bid and ask quotes. Volume and other
market-quality outcomes require genuine event-time series rather than static snapshots.

\paragraph{Matching and contract equivalence.}
Deterministic title parsing does not compare full rule text, and exact normalized keys can
hide economically important differences. Accepted matches should be checked against full
resolution rules, especially the resolution source, deadline, threshold, and cancellation
provisions. The segmented retrieval-plus-language-model procedure described above could
increase coverage while retaining human validation of the final sample.

\paragraph{Event definition and windows.}
The one-day mineral-rights window and multi-day Zelensky window are not directly comparable.
The analysis should document the key proposal, dispute, vote, resolution, and public-reporting
timestamps, then report event-time estimates under narrower and wider windows.

\paragraph{Inference.}
The reported standard errors do not account adequately for repeated observations within a
contract. Inference should cluster by matched contract, and the small number of controversy
events means the evidence should remain descriptive. Results should also be reported
unweighted and with predetermined pre-event weights because ex-post Polymarket volume may
respond to the controversy itself.

\section{Conclusion}

The project has established a useful data and matching infrastructure for studying
institutional risk in prediction markets. It has not yet established that UMA resolution
controversies caused a broad loss of price informativeness. In the matched sample,
Polymarket prices generally become more, rather than less, extreme after the two episodes,
including relative to Kalshi after the Zelensky-suit controversy. These findings are worth
reporting precisely because they discipline the theory: a general attenuation-toward-0.5
mechanism is too simple for the available data. Further work should focus on contract
validation, synchronized market data, event
timing, and defensible inference.

\begin{thebibliography}{99}

\bibitem[Rezaei Sanjabi(2026)]{fisman2026}
Rezaei Sanjabi, Rouzbeh. 2026.
``The Semantic Risk Premium.'' SSRN Working Paper 6126806.

\bibitem[Eskandari et al.(2021)]{eskandari2021}
Eskandari, Shayan, Mehdi Salehi, Wanyun Catherine Gu, and Jeremy Clark. 2021.
``SoK: Oracles from the Ground Truth to Market Manipulation.''
\textit{Proceedings of the 3rd ACM Conference on Advances in Financial Technologies}.

\bibitem[Froot and Dabora(1999)]{froot1999}
Froot, Kenneth A., and Emil M. Dabora. 1999.
``How Are Stock Prices Affected by the Location of Trade?''
\textit{Journal of Financial Economics} 53 (2): 189--216.

\bibitem[G\^arleanu and Pedersen(2011)]{garleanu2011}
G\^arleanu, Nicolae, and Lasse Heje Pedersen. 2011.
``Margin-Based Asset Pricing and Deviations from the Law of One Price.''
\textit{Review of Financial Studies} 24 (6): 1980--2022.

\bibitem[Hanson, Oprea, and Porter(2006)]{hanson2006}
Hanson, Robin, Ryan Oprea, and David Porter. 2006.
``Information Aggregation and Manipulation in an Experimental Market.''
\textit{Journal of Economic Behavior \& Organization} 60 (4): 449--459.

\bibitem[Manski(2006)]{manski2006}
Manski, Charles F. 2006.
``Interpreting the Predictions of Prediction Markets.''
\textit{Economics Letters} 91 (3): 425--429.

\bibitem[Ottaviani and S\o rensen(2007)]{ottaviani2007}
Ottaviani, Marco, and Peter Norman S\o rensen. 2007.
``Outcome Manipulation in Corporate Prediction Markets.''
\textit{Journal of the European Economic Association} 5 (2--3): 554--563.

\bibitem[Rhode and Strumpf(2008)]{rhode2008}
Rhode, Paul W., and Koleman S. Strumpf. 2008.
``Manipulating Political Stock Markets: A Field Experiment and a Century of
Observational Data.'' Working paper.

\bibitem[Shleifer and Vishny(1997)]{shleifer1997}
Shleifer, Andrei, and Robert W. Vishny. 1997.
``The Limits of Arbitrage.'' \textit{Journal of Finance} 52 (1): 35--55.

\bibitem[Wolfers and Zitzewitz(2004)]{wolfers2004}
Wolfers, Justin, and Eric Zitzewitz. 2004.
``Prediction Markets.'' \textit{Journal of Economic Perspectives} 18 (2): 107--126.

\bibitem[Wolfers and Zitzewitz(2006)]{wolfers2006}
Wolfers, Justin, and Eric Zitzewitz. 2006.
``Interpreting Prediction Market Prices as Probabilities.'' NBER Working Paper 12200.

\end{thebibliography}

\end{document}
